\section{总结与延伸}

\subsection{本讲要点回顾}

本讲介绍了 Flow Matching 作为一种全新的生成模型训练范式，其核心贡献包括：

\begin{enumerate}
    \item \textbf{连续正则化流（CNF）框架}：基于 ODE 的确定性生成过程，由流映射、向量场和概率密度路径三个核心概念构成
    \item \textbf{Flow Matching 目标函数}：简单的向量场回归损失，完全避免了 ODE 求解的计算开销
    \item \textbf{Conditional Flow Matching}：通过条件化技巧，将不可解析的边际量转化为有解析形式的条件量，同时保证训练等价性
    \item \textbf{与扩散模型的统一}：揭示了 Flow Matching 和扩散模型在数学上的等价性，以及 Flow Matching 在灵活性方面的优势
\end{enumerate}

\begin{importantbox}{关键公式速查}
\textbf{OT 路径的条件流映射}：$\phi_t(x_0 | x_1) = (1 - t)x_0 + t \, x_1$

\textbf{OT 路径的条件向量场}：$u_t(x | x_1) = x_1 - x_0$

\textbf{CFM 训练损失}：$\mathcal{L} = \mathbb{E}_{t, x_1, x_0} \left\| v_\theta(t, (1-t)x_0 + tx_1) - (x_1 - x_0) \right\|^2$

\textbf{ODE 采样}：$\frac{dx_t}{dt} = v_\theta(t, x_t), \quad x_0 \sim \mathcal{N}(0, I)$
\end{importantbox}

\subsection{拓展阅读}

\begin{itemize}
    \item \textbf{Flow Matching 原始论文}：Lipman et al., ``Flow Matching for Generative Modeling'', ICLR 2023
    \item \textbf{Conditional Flow Matching}：Tong et al., ``Improving and Generalizing Flow-Based Generative Models with Minibatch Optimal Transport'', ICML 2024
    \item \textbf{Rectified Flow}：Liu et al., ``Flow Straight and Fast: Learning to Generate and Transfer Data with Rectified Flow'', ICLR 2023
    \item \textbf{Neural ODE}：Chen et al., ``Neural Ordinary Differential Equations'', NeurIPS 2018 ——连续正则化流的开创性工作
    \item \textbf{Stochastic Interpolants}：Albergo \& Vanden-Eijnden, ``Building Normalizing Flows with Stochastic Interpolants'', ICLR 2023
    \item \textbf{Meta Flow Matching Tutorial}：\href{https://ai.meta.com/blog/flow-matching-guide-generative-ai/}{Meta AI Blog} 提供了详细的教程和代码
    \item \textbf{Fokker-Planck 方程推导}：关于向量场、概率路径和连续性方程之间关系的详细推导，可参考 Song et al., ``Score-Based Generative Modeling through Stochastic Differential Equations''（附录部分）
\end{itemize}

%% --- Fin del contenido --- %%

\end{document}
